\documentclass{article}
\usepackage[T1]{fontenc}
\usepackage{lmodern}
\usepackage{graphicx}
\usepackage{amsmath,amssymb}
\usepackage[a4paper,margin=2cm]{geometry}
\usepackage[absolute,overlay]{textpos}
\setlength{\TPHorizModule}{1mm}
\setlength{\TPVertModule}{1mm}
\usepackage[indonesian]{babel}
\usepackage{hyperref}
\setlength{\emergencystretch}{2em}
% unit: Halaman 31

\begin{document}

% === Halaman 31 (llm) ===

\begin{verbatim}
# Bungkus kunci dalam format PKCS1_OAEP untuk padding
public_key = PKCS1_OAEP.new(key.publickey())
private_key = PKCS1_OAEP.new(key)

# Tampilkan informasi kunci
print("Modulus (n):", key.n)
print("Kunci publik (e):", key.e)
print("Kunci privat (d):", key.d)
print("Bilangan prima p:", key.p)
print("Bilangan prima q:", key.q)

return public_key, private_key

def encrypt_message(message, public_key):
    #Enkripsi pesan dengan PKCS1_OAEP
    return public_key.encrypt(message.encode())

def decrypt_message(ciphertext, private_key):
    #Dekripsi pesan dengan PKCS1_OAEP
    return private_key.decrypt(ciphertext).decode()
\end{verbatim}

\end{document}
